We want to show that the functional

\begin{align*}
J(u)=\int_U |\nabla u|^2\,dx-\int_U fu\,dx
\end{align*}
attains its minimum on $H_0^1(U)$, i.e. that there exists $u_0\in H_0^1(U)$ such that
\begin{align*}
J(u_0)=\inf_{u\in H_0^1(U)}J(u).
\end{align*}
This is the standard \DEFINE{direct method in the calculus of variations}, one of the foundational techniques in Calculus of Variations.
\begin{description}

\item \textbf{Step 1: Define the infimum}

Let
\begin{align*}
m=\inf_{u\in H_0^1(U)}J(u).
\end{align*}
Choose a minimizing sequence $(u_j$\subset H_0^1$U)$ such that
\begin{align*}
J(u_j)\to m.
\end{align*}
So
\begin{align*}
\int_U |\nabla u_j|^2\,dx-\int_U fu_j\,dx\to m.
\end{align*}
We want to show $(u_j)$ has a weakly convergent subsequence.

\item \textbf{Step 2: Show the minimizing sequence is bounded}

We estimate the linear term using Cauchy–Schwarz:
\begin{align*}
\left|\int_U fu_j\,dx\right| \le \|f\|_{L^2(U)}\|u_j\|_{L^2(U)}.
\end{align*}
Using Poincaré:
\begin{align*}
\|u_j\|_{L^2(U)} \le C\|\nabla u_j\|_{L^2(U)}.
\end{align*}
Hence
\begin{align*}
\left\|\int_U fu_j\,dx\right| \le C|f\|_{L^2(U)}\,\|\nabla u_j\|_{L^2(U)}.
\end{align*}
Thus
\begin{align*}
J(u_j) \ge
\|\nabla u_j\|_{L^2}^2 - C\|f\|_{L^2}\|\nabla u_j\|_{L^2}.
\end{align*}
Let
\begin{align*}
t_j=\|\nabla u_j\|_{L^2}.
\end{align*}
Then
\begin{align*}
J(u_j)\ge t_j^2-Ct_j.
\end{align*}
Since $J(u_j)\to m$, the values $J(u_j)$ are bounded above, so $t_j^2-Ct_j$ is bounded above, which implies $t_j$ is bounded.

Therefore,
\begin{align*}
(u_j)\text{ is bounded in }H_0^1(U).
\end{align*}

\item \textbf{Step 3: Extract a weakly convergent subsequence}

Since $H_0^1(U)$ is reflexive, there exists a subsequence $still denoted (u_j)$ and some $u_0\in H_0^1(U)$ such that
\begin{align*}
u_j\rightharpoonup u_0 \qquad\text{weakly in }H_0^1(U).
\end{align*}

\item \textbf{Step 4: Use weak lower semicontinuity}

The map
\begin{align*}
u\mapsto \int_U |\nabla u|^2\,dx
\end{align*}
is weakly lower semicontinuous, so
\begin{align*}
\int_U |\nabla u_0|^2\,dx \le \liminf_{j\to\infty} \int_U |\nabla u_j|^2\,dx.
\end{align*}
Also, the linear functional
\begin{align*}
u\mapsto \int_U fu\,dx
\end{align*}
is weakly continuous, so
\begin{align*}
\int_U fu_j\,dx \to \int_U fu_0\,dx.
\end{align*}
Therefore,
\begin{align*}
J(u_0) = \int_U |\nabla u_0|^2\,dx-\int_U fu_0\,dx
\end{align*}
satisfies
\begin{align*}
J(u_0) \le \liminf_{j\to\infty} J(u_j).
\end{align*}
Since $J(u_j)\to m$,
\begin{align*}
J(u_0)\le m.
\end{align*}
But $m$ is the infimum, so
\begin{align*}
J(u_0)=m.
\end{align*}
Thus the infimum is attained.

\item \textbf{Final conclusion}

There exists $u_0\in H_0^1(U)$ such that
\begin{align*}
\boxed{
J(u_0)= \inf_{u\in H_0^1(U)} \left( \int_U |\nabla u|^2\,dx-\int_U fu\,dx \right).
}
\end{align*}

So the minimization problem has a minimizer.

\item \textbf{Why this works}

The direct method always uses:

		\begin{itemize}
        \item 1. Minimize:

Take a minimizing sequence.

        \item 2. Compactness:

Show boundedness, then extract weak convergence.

        \item 3. Lower semicontinuity:

Pass to the limit.

		\end{itemize}

This gives existence of a minimizer.

\item \textbf{PDE interpretation}

The minimizer $u_0$ is the weak solution of
\begin{align*}
-\Delta u=\frac{f}{2}
\end{align*}
(up to constants in the functional).

Thus the variational problem corresponds to an elliptic PDE.

This variational viewpoint is central in Partial Differential Equations.

---
\end{description}

If you'd like, next I can show how to derive the \DEFINE{Euler–Lagrange equation} of this minimization problem and recover the PDE satisfied by the minimizer.

